% !TEX root = ../../../main.tex
% Sources: LEC01, IMG_9110; LEC03, IMG_9128--IMG_9130.
\section{Induction and the Archimedean property}
\label{sec:analysis-i:mathematical-induction}
\label{sec:analysis-i:archimedean}

\subsection{The natural-number tools}
Mathematical induction proves a statement $P(n)$ for every $n\in\N$ by
proving $P(1)$ and proving that $P(n)$ implies $P(n+1)$.
For example, $2^n>n$ holds at $n=1$. If it holds at $n$, then
$2^{n+1}>2n\geq n+1$, which establishes the next case. The same estimate
also follows from the binomial expansion of $(1+1)^n$.

We also use the \emph{well-ordering principle}, that every nonempty subset
of $\N$ has a least element. This is equivalent to induction. To see the
direction needed here, a subset with no least element cannot contain $1$;
if it contains none of $1,\ldots,n$, it cannot contain $n+1$ either, since
that would be its least element. Induction would make the subset empty.

\subsection{Small steps eventually pass any bound}
\begin{theorem}[Archimedean property]
  \label{thm:analysis-i:archimedean}
  For every $\varepsilon>0$ and every $M\in\R$, there is an $n\in\N$
  such that $n\varepsilon>M$.
\end{theorem}
\begin{proof}
  Suppose $S=\{n\varepsilon:n\in\N\}$ were bounded above. By LUBP,
  it would have a supremum $L$. Since $L-\varepsilon<L$, there would be
  an $m\in\N$ with $m\varepsilon>L-\varepsilon$. It would follow that
  $(m+1)\varepsilon>L$, contradicting that $L$ bounds $S$. Thus $S$ is
  unbounded above, which proves the assertion.
\end{proof}
This proof uses only LUBP, so it supplies the Archimedean hypothesis used
in the preceding completeness chain. Taking the step size to be $1$ shows
that $\N$ is unbounded in $\R$. In particular, given $\delta>0$, choose
$N>1/\delta$; then $1/n<\delta$ for all $n\geq N$, proving $1/n\to0$.

\subsection{The integer part}
\begin{proposition}
  For every $x\in\R$, there is a unique $m\in\Z$ such that
  \[
    m\leq x<m+1.
  \]
  This integer is the \emph{floor} or \emph{integer part} of $x$, denoted
  by $\lfloor x\rfloor$.
\end{proposition}
\begin{proof}
  First suppose $x\geq0$. By the Archimedean property, the set of natural
  numbers greater than $x$ is nonempty. Let $n$ be its least member.
  Then $n-1\leq x<n$, so $m=n-1$ has the required property. For a general
  real $x$, choose $k\in\N$ with $x+k\geq0$. The first case gives an
  integer $r$ with $r\leq x+k<r+1$, and $m=r-k$ works.
  Finally, two distinct integers differ by at least one, so they cannot
  both satisfy the displayed inequalities. This proves uniqueness.
\end{proof}
The lecture writes the integer part as $[x]$; we use $\lfloor x\rfloor$
to distinguish it from an interval. For example, $\lfloor\pi\rfloor=3$,
whereas $\lfloor-\pi\rfloor=-4$. The floor is not truncation toward zero.

\subsection{Exercise on bounds and convergence}
An explicit sequence connects the definitions of Cauchy convergence,
infimum, and supremum. The Archimedean property supplies the indices
needed in its estimates.

% !TEX root = ../../hw02.tex
% Homework 02, Exercise 4. Shared problem statement.
\begingroup
\renewcommand{\labelenumi}{(\alph{enumi})}
\begin{exercise}
  \label{ex:hw02:decreasing-sequence}
  \solutionlink{hw02:decreasing-sequence}
  Consider the seqeunce
  \[
    x_n=\frac{1}{n^2}+\frac{2}{n^2}+\cdots+\frac{n}{n^2},
    \quad n=1,2,\ldots.
  \]
  \begin{enumerate}
    \item Show that it is monotone decreasing and bounded below.
    \item Show that it is Cauchy.
    \item Obtain the limit of $\{x_n\}$.
    \item Let $S=\{x_1,x_2,\ldots,x_n,\ldots\}$.
      Obtain both $\inf S$ and $\sup S$.
  \end{enumerate}
\end{exercise}
\endgroup

